\documentclass[../../main.tex]{subfiles}
\begin{document}


{\bf [解]}
$f(10) = 100a + 10b + c, \, g(10) = 100c + 10b + a$ で 自然数$a,b,c$ が $1 \le a,b,c \le 6$ を満たすから，題意の条件を満たすのは以下の通りである．
\begin{align}
f(10) > 452 &\iff a=5,6 \text{ 又は } \{a=4 \land 10b+c > 52\} \\
&\iff a=5,6 \text{ 又は } \text{“}a=4 \land b=6\text{”} \text{ 又は } \text{“}a=4 \land b=5 \land c \ge 3\text{”} \\
g(10) > 452 &\iff c=5,6 \text{ 又は } \text{“}c=4 \land b=6\text{”} \text{ 又は } \text{“}c=4 \land b=5 \land a \ge 3\text{”}
\end{align}
である. そこで$a,c$について以下のように排反に場合分けする.
\begin{enumerate}
  \item[$1^\circ$] $a=5,6$ かつ $c=5,6$
  \item[$2^\circ$] $a=5,6$ かつ $c=4$
  \item[$3^\circ$] $a=4$ かつ $c=5,6$
  \item[$4^\circ$] $a=4$ かつ $c=4$
\end{enumerate}

\subparagraph{$1^\circ$ の時}

 $b$ はなんでも条件は満たされるから全体の場合の数は $2 \times 6 \times 2 = 24$ 通り.

\subparagraph{$2^\circ$ の時}

 $b=5,6$ なら良く, 全体の場合の数は$2 \times 2 = 4$ 通り.

\subparagraph{$3^\circ$ の時}
 対称性から $2^\circ$ と同じく 4 通り.

 \subparagraph{$4^\circ$ の時}
  $b=5,6$ ならよく，全体の場合の数は $2$ 通り.
\casesend


全事象は $6^3$ 通りで同様に確からしいので，求める確率は
\begin{align} 
  \dfrac{24 + 4 + 4 + 2}{6^3} = \dfrac{17}{108} 
\end{align}
である．

\end{document}
